\documentclass[conference]{IEEEtran}
\IEEEoverridecommandlockouts
\usepackage{cite}
\usepackage{amsmath,amssymb,graphicx,booktabs,array,url}
% RC-20261020-v1.1; PAPER-20261020-001. Method/protocol draft; no formal results.
% Base: paper/overleaf-20261020 @ f38944669dce2cd45e124af5edc5f48a24608a00.
\title{Predictive Human Search Responsibility for\\Collaborative Multi-USV Probabilistic Search}
\author{Anonymous Authors}
\begin{document}
\maketitle
\begin{abstract}
In collaborative maritime search, taking control of one vehicle changes how much search effort remains available to the autonomous fleet. Its current position describes where that vehicle is, but does not specify where the operator intends to search or guarantee that an intended region has been observed. This paper formulates human search responsibility as a planning state separate from target belief. An explicit SEARCH declaration activates a responsibility anchor; navigation and control interruptions revoke it. A belief-weighted allocation retains every searchable cell while adapting autonomous centers around the anchor. This formulation defines a controlled comparison between the human-controlled USV's current position and a valid predicted future position, with reactive fallback when prediction is unavailable. The implemented simulation infrastructure supports declaration, revocation, and operation without a prediction model. Comparative evaluation is specified using a shared sensing and planning pipeline and an independently assessed sustained-proximity encounter endpoint. The present evidence supports the engineering foundation; whether future-position information improves collaboration remains to be tested.
\end{abstract}

\begin{IEEEkeywords}
human--robot collaboration, multi-robot search, unmanned surface vehicles, responsibility allocation, motion prediction
\end{IEEEkeywords}

\section{Introduction}
A team searching for an unknown target must decide both where to search and how to distribute its effort. In probabilistic search, sensor observations update an estimate of target location, and planning uses that estimate to select subsequent actions~\cite{b1,b7}. Human intervention introduces a second source of uncertainty: the autonomous fleet may know which vehicle an operator controls without knowing where that vehicle will contribute search effort. A useful collaboration mechanism must account for that contribution without confusing the operator's actions with evidence about the target.

Consider an operator who takes control of one USV and begins searching toward another part of the workspace. Representing its contribution only by its current position supplies a reactive allocation reference. A future-position estimate could supply an earlier reference, but the estimate is meaningful only while the operator is actually engaged in search. Takeover may instead be used to transit, inspect the interface, or reposition the vehicle. Moreover, an intended visit does not establish that the region has been sensed. These distinctions motivate separate representations of execution authority, search commitment, and target evidence.

Anticipating human search is not a new idea. Previous systems predict the next human search location to guide complementary robot action~\cite{Xie2011Intention}, incorporate anticipated human-searcher paths into robot trajectory planning~\cite{Heintzman2021Anticipatory}, and adapt autonomous search to the behavior of a teleoperated teammate~\cite{KrzysiakButail2025Movement,bb2025}. These precedents motivate a specific question here: \emph{during an explicitly declared SEARCH task, what decision value does the predicted future position of the human-controlled USV provide beyond its current position within the same responsibility-allocation mechanism?}

We formulate \emph{human search responsibility} as an active planning commitment represented by a spatial anchor. It is activated by an explicit SEARCH declaration, rather than inferred from takeover alone. NAVIGATE, cancellation, release, and vehicle switching withdraw obsolete responsibility. The anchor represents the operator's contribution through the controlled USV; it is not the operator's physical position or a measured target location. All searchable cells remain in the allocation domain, including cells associated with that anchor. Only actual sensor reports update the target belief.

The completed work establishes this responsibility-state formulation, its connection to fixed-center belief-weighted allocation, and an evaluation protocol that distinguishes current-position information (B2) from predicted-position information (B3). Shared inference, allocation, path-cost assignment, and execution define the comparison; the anchor information is the intended difference. The simulation implementation has undergone bounded engineering checks of explicit SEARCH, revocation, and current-position fallback without a learned model. These checks establish an executable foundation, not the performance of a trained predictor or the benefit of prediction.

Accordingly, this manuscript reports a method and evaluation design rather than comparative search results. The primary outcome is time to an independently evaluated sustained-proximity encounter, distinct from sensor detection or visual confirmation. Future experimental findings can establish whether prediction improves this outcome; B2--B3 alone cannot establish that the proposed allocation mechanism is superior to every alternative use of a human forecast.

\section{Related Work}
\subsection{From Target Belief to Spatial Search Allocation}
Probabilistic search makes target uncertainty explicit. Bourgault \emph{et al.} study coordinated Bayesian search~\cite{b1}, while Hollinger \emph{et al.} address multirobot search for a moving target~\cite{Hollinger2009Efficient}. Dames combines distributed search and tracking with a PHD representation that accommodates an unknown number of targets, finite sensor fields of view, missed detections, and clutter~\cite{b7}. The present study instead assumes one stationary target and a normalized categorical belief. Its observation model supplies a common task state for comparing collaboration mechanisms; it is not a new target filter.

Spatial decomposition provides a related allocation foundation. Locational coverage objectives and centroidal Voronoi/Lloyd methods place sensing resources according to a spatial density~\cite{Cortes2004Coverage}. Local Voronoi task allocation~\cite{Fu2009LocalVoronoi} and coverage optimization with static and mobile sensors~\cite{Mahboubi2013StaticMobile} further establish that spatial resources and their geometry can structure mobile-agent effort. Our fixed-center allocation builds on these ideas. The human anchor, however, represents a prospective search contribution whose execution remains uncertain, rather than proof that its associated region has already been observed. Neither the partitioning algorithm nor the presence of a fixed generator is claimed as an independent contribution.

\subsection{Anticipating Human Search to Obtain Complementary Action}
Human-search prediction has been used at several decision levels. Xie \emph{et al.} learn hidden-Markov models of human search and predict the next observation location~\cite{Xie2011Intention}. The prediction supplies inhibition information to a temporary planning PDF that guides complementary robot actions. Their framework distinguishes this planning PDF from the target PDF updated by robot observations, and also learns object priors from human search-activity density. This distinction matters: their use of human information cannot be reduced to treating a prediction as a permanent target observation.

Heintzman \emph{et al.} incorporate anticipated ground-searcher paths, a dynamic lost-person model, and UAV sensing into posterior-risk trajectory optimization~\cite{Heintzman2021Anticipatory}. At the search-region level, Ishii \emph{et al.} estimate areas searched by a human, predict future search areas, and use dialog to resolve ambiguity~\cite{Ishii2023Enhancing}. Lim Hong Da \emph{et al.} develop this space-sharing setting through on-demand bidirectional interaction~\cite{LimMiuraHayashi2025SpaceSharing}. Together, these works establish complementary search based on forecasts of human locations, paths, or regions. Predicting where a person will search and avoiding prospective duplication are therefore precedents, not claims of novelty here.

They also motivate care in interpreting responsibility. Xie \emph{et al.} report cases in which a human misses an object and inhibition discourages the robot from searching the same location~\cite{Xie2011Intention}. This is a limitation of the combined estimation and cooperation scheme, not isolated evidence about prediction error. Our formulation retains target belief and the complete search domain, but this alone does not guarantee timely revisitation. The relevant distinction to evaluate is how a forecast changes allocation and realized sensing, rather than whether a system assigns a new name to anticipated human effort.

\subsection{Adapting Search to a Teleoperated Teammate}
The human need not be a separate ground searcher. Krzysiak and Butail formulate information-based control in which autonomous robots balance target search against assistance near a human-teleoperated reference robot~\cite{KrzysiakButail2022Information}. Their later work infers prior knowledge from movement cues and adapts mutual-information blending; its simulation uses human trajectories recorded during experiments~\cite{KrzysiakButail2025Movement}. Bhattacharya and Butail similarly infer operator knowledge and situational awareness from teleoperated UAV motion and adapt swarm separation and attraction~\cite{bb2025}. These are close precedents for using the behavior of a controlled fleet member to organize autonomous search. Same-team teleoperation, behavioral adaptation, and recorded-trajectory replay are not new elements of the present study.

Other collaborative frameworks expose different places for human information to enter planning. Dalmasso \emph{et al.} use Social Reward Sources and multi-agent MCTS for collaborative navigation and search~\cite{b5}; Talebpour and Martinoli incorporate predicted human motion and uncertainty into risk-based task bids~\cite{TalebpourMartinoli2018IROS}. Our comparison concerns a particular allocation input: the human-controlled USV's current or predicted position is held fixed as a generator while the remaining autonomous generators are updated over the full belief-weighted domain. An explicit SEARCH state specifies when that input represents a search commitment.

This scope yields a testable information comparison without asserting that all previous systems lack equivalent commitment, revocation, or fallback behavior. It also separates two questions. B2--B3 tests future-position information within the same allocation operator. A comparison against direct predictive inhibition would instead test the choice of allocation mechanism; it remains an auxiliary replay candidate, not a required main condition or an advantage already demonstrated by this work.

\section{Problem Formulation}
\subsection{Environment and Inference}
A known planar obstacle map defines searchable free cells $\mathcal X=\{x_i\}_{i=1}^N$. Exactly one stationary target occupies an unknown free cell. A fleet of $n$ USVs has positions $p_r(t)$; the VRX platform is configured for ten WAM-Vs. Hidden target truth has two legitimate uses: generating synthetic sensor reports and independently evaluating target encounters. Neither the planner nor the target estimator receives the true target position, its target-specific distance, or truth-derived admission and abstention reasons. Public distances to hypothetical target cells remain valid model inputs.

\subsection{Effective Sensor Opportunities}
Let $e_k$ identify an accepted sensor opportunity and $G_k$ contain its acquisition geometry and public input-validity information. Each opportunity has one identifier scoped to the mission, sensor, and opportunity index; retransmission does not create new evidence. Under the current engineering adapter, a valid opportunity yields a single binary report $z_k\in\{0,1\}$ for the sensor as a whole. The grid cells are alternative hypotheses about $X$, not separate independent reports from that one opportunity. ABSTAIN denotes absent evidence and is not a third detection outcome.

We define a sensor-evidence model $\mathcal M$ with observation history $\mathcal H_{k-1}=Z_{1:k-1}$. Acquisition geometries $G_{1:k}$ are supplied design/control inputs, denoted after a semicolon. Within this model, their selection contributes no extra likelihood about $X$:
\begin{align}
 b_i^k&=P_{\mathcal M}(X=x_i\mid Z_{1:k};G_{1:k}),\\
 L_{i,k}&=p_{\mathcal M}(z_k\mid X=x_i,\mathcal H_{k-1};G_{1:k}),\\
 b_i^k&=\frac{b_i^{k-1}L_{i,k}}{\sum_j b_j^{k-1}L_{j,k}},\quad \sum_i b_i^k=1.
\label{eq:bayes}
\end{align}
The recursion intentionally omits an action/geometry-selection likelihood. It uses $z_k$ once to update $b^{k-1}$ to $b^k$ under $\mathcal M$. This is a modeling restriction, not evidence that human actions are independent of the target. If an operator acts on additional target knowledge, $b$ is not claimed to equal a posterior exploiting all available human-behavior information. The normalizing denominator must be positive; an impossible report or zero denominator is a model or numerical fault, not grounds for adding arbitrary probability mass.

In the assumed binary engineering model, $p_{i,k}=P_{\mathcal M}(z_k=1\mid X=x_i;G_k)$ and
\begin{equation}
 L_{i,k}=p_{i,k}^{z_k}(1-p_{i,k})^{1-z_k}.
\label{eq:binary}
\end{equation}
The main comparison uses this idealized binary model with zero false alarms and no reported target coordinate. Its positive likelihood is $p_{i,k}$, not a localized measurement of the true target position; its negative likelihood is $1-p_{i,k}$. A positive report does not establish mission success. Extending the adapter to false alarms or measured positions requires a new observation likelihood, not reinterpretation of this binary output.

The existing engineering simulator uses a hard horizontal field of view and logistic range envelope:
\begin{equation}
 p_{i,k}=v_{i,k}\frac{p_{\max}}{1+\exp[(d_{i,k}-r_{50})/s]},
\label{eq:pd}
\end{equation}
where $v_{i,k}$ gates field of view, configured range support, and modeled visibility. The parameters and effective-opportunity spacing are assumed simulation settings; they are not estimated camera performance. Given $X$ and public geometry, independent Bernoulli draws define the main simulation profile, not a claim of independence among real cameras or across real observation times. Its idealized visibility convention sets $p_{i,k}=0$ for both occluded and geometrically unknown cells. For such cells a negative report has neutral likelihood one; a whole-sensor positive report has likelihood zero. This convention does not guarantee a conservative posterior or a lower bound on real detection under unknown visibility. An entirely absent opportunity instead contributes no likelihood anywhere. Normalization may increase posterior mass outside the negatively observed support without adding mass manually.

Public pose, timing, and map validity determine opportunity admission before consulting the hidden target. In simulation, the sampler computes the same grid probability vector as the estimator, indexes it at the hidden target cell, and samples the binary outcome. Truth may determine the generated outcome but must not determine whether an opportunity is published, its validity, or its public abstention reason. Invalid truth/state-space configuration is a simulation fault, not an ordinary location-dependent abstention. This is a generative simulation boundary; it does not claim image-based perception.

Multiple reports may be processed sequentially. The chain-rule joint likelihood is a product of history-conditioned factors. Replacing these with the marginal binary factors in~\eqref{eq:binary} assumes conditional independence given the target and modeled geometry. Under that explicit simulation assumption, sequential normalization and one batch product are equivalent; shared scene effects or reused detector windows can invalidate the factorization. Sensor timestamps within one opportunity window need not be simultaneous. Unique identifiers prevent duplicate use but do not establish statistical independence.

For a negative-only history under this model, the diagnostics
\begin{equation}
 S_i^k=\prod_{u=1}^k(1-p_{i,u}),\qquad R^-_k=\sum_i b_i^0S_i^k
\end{equation}
are unnormalized no-positive-report compatibility and its prior-weighted total. Then $b_i^k=b_i^0S_i^k/R^-_k$. If positive reports also update belief while the implementation retains these negative-history diagnostics, that identity no longer represents the mixed-history posterior. These quantities are not calibrated absolute mission risk in this study.

Commands, gaze, joystick inputs, declarations, and predictions do not directly enter~\eqref{eq:bayes}. Legitimate sensing continues in SEARCH and NAVIGATE; geometric visitation alone does not clear a cell. Camera, infrared, sonar, and LiDAR target adapters are not interchangeable. Four-modality fusion is outside the main experiment. The numerical observation parameters, support, opportunity spacing, and input-validity settings must be exported from the final implementation into a consistent manifest before formal testing.

\subsection{Explicit Human Search Responsibility}
The operator directly controls at most one vehicle $h(t)$. Directly controlled and active command-locked vehicles are excluded from autonomous reassignment. An execution goal and a search anchor are separate: retaining a NAVIGATE command does not require using its destination as a search center.

Let $a_h(t)$ indicate active explicit SEARCH. Takeover initializes $a_h=0$. A discrete declaration activates it, and the state persists between actions. NAVIGATE, cancellation, release, control-segment termination, or switching vehicles revokes it. A new vehicle requires a new declaration. The declaration answers whether the operator searches, not where the vehicle will be; no advance destination is required.

\section{Predictive Responsibility Redistribution}
\subsection{Current and Future Anchors}
During active SEARCH, the B2 and B3 anchors are
\begin{align}
 c_h^{B2}(t)&=p_h(t),\\
 c_h^{B3}(t)&=\begin{cases}\hat p_h(t+H),&v_h(t)=1,\\p_h(t),&v_h(t)=0,\end{cases}
\label{eq:anchor}
\end{align}
where $H$ is prediction horizon and $v_h$ denotes validity of the current vehicle and control/task segment. Validity is a predictor-input gate, separate from the active SEARCH state. Out-of-domain or infeasible endpoints must fail the gate rather than reserve an unreachable responsibility point; the admissibility test is to be frozen. Without active SEARCH, neither condition inserts a human search anchor. Missing prediction must not block fleet planning. Horizon, validity criteria, refresh rate, and prediction-driven replanning thresholds remain to be determined independently of final mission outcomes.

The prediction target is the current human-controlled USV position $p_h(t+H)$, not the operator's body position, a selected navigation goal, target-object position, or an intention label. RF regression is the first development candidate; it is neither a validated model nor the automatically selected final predictor. The development protocol calls for actual operator inputs and trajectories recorded in VRX, distinct from scripted fixtures and from physical sea trials. Developer-only recordings support a developer pilot, not population-level generalization.

\subsection{Causal Prediction Interface}
At a prediction request with simulation event time $t$, features must already have arrived and have source timestamps no later than $t$. A source timestamp alone does not make a late-arriving record available retrospectively. Request identity and the available source-record snapshot must be retained. The forecast records its base time $t$, target time $t+H$, generation time, mission/vehicle/segment identities, and model version; generation delay must not shift the target time. A forecast whose target time has passed is rejected.

The initial feature candidate is body-frame forward and lateral velocity, yaw rate, current forward and turning commands, and the ages of the reference pose and command. Velocities may use documented causal differences between valid past poses from the same segment. Source ages use acquisition times; issued operator commands remain distinct from executed thrust. Missing quantities are not silently zero-filled. Gaze, speech, navigation goals, target belief, hidden target state/distance, future commands, and final outcomes are excluded from this RF profile. Coordinates and heading define the output anchor/rotation, rather than default absolute-map features. Source identity must be attached at acquisition and checked on receipt; downstream callbacks cannot manufacture provenance by relabeling an old sample with the current segment.

\subsection{Offline Training Labels}
Training inputs and offline future labels remain separate. A label requires the same mission, vehicle, and control segment with uninterrupted explicit SEARCH throughout $[t,t+H]$. Cancellation, release, switching, NAVIGATE, or mission termination before the target time makes the label unavailable, with its reason retained. Segment-end eligibility uses the owner's event order and never gates online prediction using knowledge of the future. Where no exact target-time pose exists, offline label generation may interpolate only between valid bracketing poses in the same segment, under a maximum gap frozen before processing. Without both endpoints or a justified gap limit, no label is created; extrapolation and unbounded forward filling are disallowed. This permission concerns prediction labels only: encounter evaluation retains its no-interpolation rule.

A label may be represented as body-frame future displacement relative to the latest legal reference pose at $t_{\rm ref}\leq t$. Its timestamp, frame, heading, and transform version are retained; reconstruction still predicts world position $p_h(t+H)$. The reference pose is not assumed to have been acquired exactly at the request time. Numerical $H$, age/gap limits, feature preprocessing, and RF hyperparameters remain to be selected on development evidence and frozen before formal evaluation.

\subsection{Fixed-Center Allocation and Execution}
With fixed anchors $\mathcal C_F$ and autonomous centers $\mathcal C_A$, the shared surrogate is
\begin{equation}
 \Phi=\sum_{i=1}^N b_i\min_{c\in\mathcal C_A\cup\mathcal C_F}\|x_i-c\|_2^2.
\label{eq:cluster}
\end{equation}
Fixed anchors remain fixed while autonomous centers update. All searchable free cells remain in the responsibility domain with deterministic tie-breaking. The nonconvex objective is approximated with fixed-generator weighted Lloyd updates, not solved with a claim of global optimality. For an autonomous cell set $\mathcal V_j$ with positive mass, its centroid is
\begin{equation}
 \bar c_j=\frac{\sum_{i\in\mathcal V_j} b_i x_i}{\sum_{i\in\mathcal V_j}b_i}.
\end{equation}
The candidate center must be admissible in free space. Empty or zero-mass clusters require a deterministic policy; their centroid formula is undefined and may not be patched by injecting probability mass. Initialization, tie-breaking, projection/reseeding, stopping tolerance, and iteration cap must match across B2 and B3 and be recorded in the final profile. Zero-weight cells contribute zero to~\eqref{eq:cluster} but still receive geometric responsibility labels. No residual region is manually deleted and no artificial probability mass is introduced to repair allocation behavior. Planned responsibility does not guarantee actual sensing or create navigation exclusions.

For each eligible autonomous vehicle--center pair, A* first computes an obstacle-aware path $\mathcal P_{rj}$. Path-dependent costs then define assignment, for example
\begin{equation}
 \min_\pi\sum_{r\in\mathcal R_A} L(\mathcal P_{r,\pi(r)}).
\end{equation}
Hungarian assignment follows cost construction. Unreachable pairs are inadmissible, not zero-cost assignments. If no complete feasible matching exists, the candidate allocation must not be committed as executable; the fallback behavior and its logging require explicit implementation acceptance. When no autonomous vehicle is eligible, no autonomous matching is performed. Path length is a candidate baseline to freeze in the experiment manifest; estimated travel time is a separate alternative. Both conditions share region coverage and tracking. A* provides geometric paths, not a guarantee of WAM-V dynamic feasibility.

\subsection{Interruption and Fallback}
Predictions and declarations must be associated with the current mission, vehicle, and control/task segment. Ending a segment invalidates its information. Delayed messages cannot reactivate obsolete responsibility; duplicate declarations should be idempotent. Existing planner-owned revisions should be reused where adequate. Recorder observes these states rather than becoming their authority.

Task and authority changes cause auditable planning updates. Prediction thresholds require validation. A stale-anchor ablation can test the interruption mechanism, but stale responsibility is not permissible normal behavior. On a new mission, controlled vehicle, or control/task segment, any old predictor window is invalidated. Until a sufficient new causal window exists, active SEARCH uses the current-position anchor. Loss of SEARCH removes the anchor entirely, rather than preserving it through fallback. Segment validity transitions and prediction changes must be distinguished in the replanning log. These are method requirements. Current-position fallback without a learned model has passed bounded engineering checks; valid learned forecasts and their end-to-end use remain unvalidated. General negotiation, proactive clarification, and new motion-planning algorithms remain outside the current contribution.

\section{Evaluation Protocol}
\subsection{Independent Sustained-Proximity Criterion}
An independent evaluator uses target truth $x^\star$. Radius $\rho_E$ and dwell $\tau_E$ define an ideal continuous-proximity criterion
\begin{equation}
\begin{split}
 T_E^{\rm cont}=\inf\{t\geq\tau_E:\;&\exists r,\ \|p_r(u)-x^\star\|_2\leq\rho_E\\
 &\text{for every }u\in[t-\tau_E,t]\}.
\end{split}
\label{eq:event}
\end{equation}
Time is measured from mission start in simulation time. Each vehicle has its own encounter timer; leaving resets it. Visits and vehicles cannot pool time. Autonomous, manual, and NAVIGATE vehicles are evaluated identically. No heading, field-of-view, or probability condition is added. The endpoint deliberately measures sustained proximity rather than whether the assumed sensor produced a positive report.

The operational endpoint is \emph{sampled sustained proximity}, denoted by $T_E$. At each valid simulation-time pose sample, the evaluator updates a separate qualifying interval for each vehicle. All consecutive samples in an interval must satisfy $d\leq\rho_E$. When its elapsed duration first reaches $\tau_E$, encounter occurs at that current sample time. There is no interpolation, backdating, or use of future samples to infer an earlier success. Leaving, an invalid pose, or a gap longer than the allowed interval resets the timer; resumption starts a new interval. Duplicate timestamps add no duration. Clock reversal or a mission-epoch change invalidates the old timing state. Sample cadence and the maximum permitted gap must be recorded and frozen. Finite samples do not prove uninterrupted physical residence between them, so $T_E$ is distinguished from $T_E^{\rm cont}$ in~\eqref{eq:event}.

This endpoint is proximity, not visual detection, confirmation, or rescue. Hidden target truth is used inside the synthetic observation generator and the independent encounter evaluator. Their interfaces expose only legitimate reports or encounter events/times for recording and termination, never hidden target coordinates or target-specific distances to planning or inference.

At a common budget $B$, the intended observations are $Y=\min(T_E,B)$ and $\delta=\mathbf1[T_E\leq B]$. Radius, dwell, and the formal budget remain symbolic until their numerical settings are frozen. Encounter is the sole normal success event in the main experiment. A random POSITIVE, legacy $\Gamma$/FOUND, or visual confirmation cannot independently declare success or normally stop search in place of encounter. Existing terminal messages may be reused only with an unambiguous encounter reason and one terminal-state owner.

The administrative deadline is $t_{\rm start}+B$ regardless of termination-message delay. A valid encounter at $T_E=B$ counts as success and takes precedence over natural budget expiration. The terminal owner must arbitrate same-step encounter and budget candidates by their simulation event times before publishing one irreversible terminal state; callback arrival order cannot extend $B$ or justify a later terminal rewrite. Safety-fault stops remain separate and are not overridden by this natural success/budget priority. Administrative STOP, invalid timing, and crashes require separate disposition rather than automatic ordinary budget censoring. These are protocol requirements, not claims of completed endpoint implementation.

\subsection{Conditions and Study Design}
The primary comparison is B2 versus B3, sharing the declaration interface, sensing model, clustering, assignment, controller, and budget. B0 all-autonomous and B1 human-control without a responsibility anchor are supporting references. Offline B4 uses actual future human motion as a perfect-information reference for the same mechanism, not a deployable predictor or guaranteed performance upper bound.

Paired replay holds human trajectory and declaration/cancellation events fixed within scenarios. Map, target placement, initial states, and sensing configuration are matched. Equal seeds alone do not imply identical observations when autonomous trajectories diverge; the coupling policy must be recorded. B2/B3 share the observation update rule and initial prior, while their realized reports and posteriors can differ after their routes diverge. Replay isolates the mechanism under fixed human behavior. A later counterbalanced live study assesses interaction effects and cannot assume identical trajectories across conditions. Participant collection follows applicable institutional reporting and approval requirements.

\subsection{Predictor Development and Evaluation Protocol}
Raw sessions and trajectories are grouped before generating samples. Every window, derived copy, and replay from one recording stays in the same partition. Cross-operator claims require complete participant holdout, including all of that participant's sessions; segment boundaries cannot justify splitting one session across training and testing. With insufficient independent participants, evaluation is limited to explicitly labeled within-operator, held-out-session development. Fitted preprocessing, feature selection, $H$, thresholds, and RF tuning use development data only. Row-wise random splits and out-of-bag scores do not replace grouped evaluation. The final predictor test and formal collaboration evaluation use development-unused groups; model and processing changes after observing those results require a new untouched evaluation set.

RF, current-position hold, and constant-velocity extrapolation are compared at the same request times and target-time labels using the same available state. Constant velocity accounts for the actual interval from the reference-pose timestamp to $t+H$. Report planar endpoint error in meters on common labelable samples, first by session and then with participant-aware aggregation. Also report all input attempts, label coverage/unavailability, legal output rate, fallback frequency and reasons, and the deployed output's error after fallback. If the current pose is invalid too, record no output rather than zero error. Conditional accuracy on continuous, labelable SEARCH intervals is not accuracy over all decisions. Selecting a shorter $H$ for lower error alone does not establish useful coordination lead time or responsibility changes. Better prediction does not imply better encounter RMST. Sample sizes, split proportions, confidence procedures, and final numerical settings remain open.

Predictor selection and parameter validation use held-out development data, not final mission improvement. Formal participant/scenario structure and the analysis protocol remain to be finalized before evaluation.

\subsection{Primary and Supporting Outcomes}
The primary estimand is
\begin{align}
 \mathrm{RMST}_c(B)&=\int_0^B P(T_E^c>t)\,dt,\\
 \Delta_E&=\mathrm{RMST}_{B3}(B)-\mathrm{RMST}_{B2}(B).
\end{align}
Negative $\Delta_E$ supports faster encounters under this protocol. Budget-limited failures are retained. For valid trials with only common-budget administrative censoring, $\mathrm{RMST}_c(B)=\mathbb E[\min(T_E^c,B)]$; thus failures contribute $B$ instead of disappearing from a success-only mean. This identity does not justify treating a crash or administrative early STOP as an ordinary budget-limited failure. In paired replay, a complete scenario, human trajectory, and declaration sequence define a paired comparison unit. Random repeats under the same trace are nested replicates, not independent operators or trajectories. Multiple traces from one participant require participant-aware clustering. Frames, opportunities, and overlapping predictor windows are not independent experimental units. Paired bootstrap intervals remain a candidate; sample size and the exact analysis must be predeclared. Infrastructure exclusions or reruns follow condition-independent rules and retain all records; method-caused crashes cannot be silently removed.

Supporting evidence includes redistribution lead time, prediction error, fallback frequency, encounter success by budget, travel, and planning cost. A candidate overlap measure uses the modeled geometric supports of accepted sensing opportunities. It must specify cross-vehicle versus same-vehicle repetition, opportunity windows, denominator, and comparability after early termination. Neither planned responsibility overlap nor modeled support overlap is automatically real detection redundancy. No unverified overlap or hazard metric is frozen here. Optional workload and declaration-burden measures address interaction cost. Faster encounter alone does not establish reduced redundancy or improved operator experience.

\section{Implementation Status and Limitations}
The system uses ROS~2 Humble, Gazebo, and VRX with ten WAM-Vs. Tests exercised the shared sensor-event interface, one-use report identifiers, negative updates, and neutral abstention. For fixed public poses, maps, and times, admission and validity do not depend on hidden target truth. A separate run exercised budget expiry, safe stopping, and recorder closure. These tests support engineering behavior under the assumed model. They do not calibrate sensor performance or establish real-world independence.

An API-level final-configuration run also exercised explicit SEARCH activation of a B2 current-position anchor and its cancellation, navigation, and release transitions within the simulation clock domain. This is bounded execution evidence, not verification of physical interface clicks, every possible race or the independent encounter evaluator. This responsibility run has passed engineering review within that API-level scope. Stage-specific evidence is retained separately from formal mission comparisons.

A final-configuration API-level run supports operation without an RF model: active SEARCH used a current-position fallback, autonomous planning continued briefly, and cancellation, navigation, and release revoked the anchor. This bounded fallback behavior has passed engineering review; it does not establish the source validity or accuracy of real RF forecasts. Encounter-based evaluation remains incomplete. Four-modality fusion and real-perception confirmation are outside the current main experiment. Engineering budget expiry is not yet an accepted encounter-censoring contract. General communication planning, prediction-risk recovery, and targeted overlap-support dispatch remain future work.

Operator-facing controls, button mappings, existing device links, acquisition-time provenance, session export, and parameterized training entry points are preparation requirements. API tests or labeled engineering fixtures do not establish actual UI/hardware validation, volunteer data collection, or completed RF training. Core operation must remain possible without an RF file or an optional gaze/speech device; optional device failures must be visible and recorded. Physical operator records and identity mappings remain locally controlled, with only schemas, versions, de-identified summaries, and necessary hashes shared publicly.

No formal B2/B3 outcome is reported for the sampled endpoint $T_E$. Required evidence includes predictor validation, responsibility/interruption execution, encounter timing/censoring, paired outcomes, and mechanism measurements. Completed RF validation, calibrated perception, gaze validation, or participant studies are not claimed.

The point anchor is a proxy, not a commitment to cover its entire induced region. Retaining belief and the full search domain does not guarantee prompt autonomous revisitation if human search is missed or a commitment is unfulfilled; the allocation may still delay that visit. Such robustness is not established by evidence--responsibility separation alone. Inaccurate prediction may harm allocation; frequent fallback may make B3 resemble B2 and must be reported. Tracking quality can confound comparisons. Assumed sensing and proximity-based success support a bounded simulation claim, not real-camera detection or field deployment.

\section{Results Pending}
No formal mission or participant results are available. After model and endpoint freezing, this section will report the paired B2--B3 encounter RMST estimate with uncertainty, encounter success, actual sensing overlap, prediction/fallback behavior, planning cost, and failures. Engineering smoke records will not be pooled with the formal comparison. Unsupported mechanism or workload improvements will not be inferred from a faster encounter alone.

\section{Conclusion}
This draft focuses on the value of predicting explicitly declared human search responsibility. B2 and B3 share inference and planning, differing in valid responsibility anchors. Interruption revokes obsolete responsibility. An independent sustained-proximity endpoint provides a simple simulation comparison. Benefits remain hypotheses pending engineering acceptance and formal experiments.
\begin{thebibliography}{99}
\bibitem{b1}
F. Bourgault, A. Goktogan, T. Furukawa, and H. F. Durrant-Whyte, ``Coordinated search for a lost target in a Bayesian world,'' \emph{Advanced Robotics}, vol. 18, no. 10, pp. 979--1000, 2004, doi: 10.1163/1568553042674707.

\bibitem{b7}
P. M. Dames, ``Distributed multi-target search and tracking using the PHD filter,'' \emph{Autonomous Robots}, vol. 44, pp. 673--689, 2020, doi: 10.1007/s10514-019-09840-9.

\bibitem{Xie2011Intention}
D. Xie, Y. Lin, R. A. Grupen, and A. R. Hanson, ``Intention-based coordination and interface design for human-robot cooperative search,'' in \emph{Proc. IROS}, pp. 3349--3354, 2011, doi: 10.1109/IROS.2011.6094688.

\bibitem{Heintzman2021Anticipatory}
L. Heintzman, A. Hashimoto, N. Abaid, and R. K. Williams, ``Anticipatory planning and dynamic lost person models for human-robot search and rescue,'' in \emph{Proc. ICRA}, pp. 8252--8258, 2021, doi: 10.1109/ICRA48506.2021.9562070.

\bibitem{KrzysiakButail2025Movement}
R. Krzysiak and S. Butail, ``Human prior knowledge estimation from movement cues for information-based control of mobile robots during search,'' \emph{ACM Trans. Human-Robot Interaction}, vol. 14, no. 2, art. 27, 2025, doi: 10.1145/3706065.

\bibitem{bb2025}
A. Bhattacharya and S. Butail, ``Robot movement based inference of teleoperator state for multi-robot search,'' \emph{Complex \& Intelligent Systems}, vol. 11, art. 443, 2025, doi: 10.1007/s40747-025-02055-9.

\bibitem{Hollinger2009Efficient}
G. A. Hollinger, S. Singh, J. Djugash, and A. Kehagias, ``Efficient multi-robot search for a moving target,'' \emph{Int. J. Robotics Research}, vol. 28, no. 2, pp. 201--219, 2009, doi: 10.1177/0278364908099853.

\bibitem{Cortes2004Coverage}
J. Cort\'es, S. Mart\'inez, T. Karatas, and F. Bullo, ``Coverage control for mobile sensing networks,'' \emph{IEEE Trans. Robotics and Automation}, vol. 20, no. 2, pp. 243--255, 2004, doi: 10.1109/TRA.2004.824698.

\bibitem{Fu2009LocalVoronoi}
J. G. M. Fu, T. Bandyopadhyay, and M. H. Ang, ``Local Voronoi decomposition for multi-agent task allocation,'' in \emph{Proc. ICRA}, pp. 1935--1940, 2009, doi: 10.1109/ROBOT.2009.5152829.

\bibitem{Mahboubi2013StaticMobile}
H. Mahboubi, J. Habibi, A. G. Aghdam, and K. Sayrafian-Pour, ``Distributed coverage optimization in a network of static and mobile sensors,'' in \emph{Proc. ACC}, 2013, doi: 10.1109/ACC.2013.6580919.

\bibitem{Ishii2023Enhancing}
T. Ishii, J. Miura, and K. Hayashi, ``Enhancing human-robot collaborative object search through human behavior observation and dialog,'' in \emph{Proc. ICCVW}, pp. 1841--1848, 2023, doi: 10.1109/ICCVW60793.2023.00199.

\bibitem{LimMiuraHayashi2025SpaceSharing}
N. Lim Hong Da, J. Miura, and K. Hayashi, ``Enhancing human-robot collaborative search through efficient space sharing with on-demand bidirectional interaction,'' in \emph{Computer Vision--ECCV 2024 Workshops}, pp. 20--35, 2025, doi: 10.1007/978-3-031-92591-7\_2.

\bibitem{KrzysiakButail2022Information}
R. Krzysiak and S. Butail, ``Information-based control of robots in search-and-rescue missions with human prior knowledge,'' \emph{IEEE Trans. Human-Machine Systems}, vol. 52, no. 1, pp. 52--63, 2022, doi: 10.1109/THMS.2021.3113642.

\bibitem{b5}
M. Dalmasso, J. E. Dom\'inguez-Vidal, I. J. Torres-Rodr\'iguez, P. Jim\'enez, A. Garrell, and A. Sanfeliu, ``Shared task representation for human--robot collaborative navigation: The collaborative search case,'' \emph{Int. J. Social Robotics}, vol. 16, pp. 145--171, 2024, doi: 10.1007/s12369-023-01067-0.

\bibitem{TalebpourMartinoli2018IROS}
Z. Talebpour and A. Martinoli, ``Risk-based human-aware multi-robot coordination in dynamic environments shared with humans,'' in \emph{Proc. IROS}, pp. 3365--3372, 2018, doi: 10.1109/IROS.2018.8593586.
\end{thebibliography}
\end{document}
